\section{Conclusion}\label{sec:conclusion}
This paper proposed \mname{}, a multi-UAV replanner that plans against a chance constraint on the finish time of
every task. Because disturbances persist over a leg, the spread of the finish-time error grows linearly with the lead
time, and with an affine error model the chance constraint reduces exactly to an effective deadline per task, so fast
heuristics are reused unchanged. A fallback keeps the tasks that the constraint would give up. On three independent
sets of held-out seeds of pre-registered studies in a kinematic simulator, the chance-constrained model served more
reward than the nominal model and a constant time buffer, and on the last set the replanner served
\val{QURelCCConv}\% more than a conventional event-triggered planner. A tuned proportional buffer and moment
propagation along the route reached the same utility, which identifies the margin that grows with the lead time,
together with the fallback, as the source of the gain in utility. On a further set of seeds, the replanner served
\val{WURelCACBBA}\% more than CBBA, and CBBA on the proposed model came close to it, so the gain lies in the planning
model and not in the search.

The probability of meeting a deadline, which a buffer lacks, lets a replanner price lateness. The cost-aware fallback
attempts a given-up task only if this probability outweighs the cost of a late arrival, and the closed form turns the
rule into another effective deadline without a parameter to tune. With a late arrival that costs one task reward, it
served \val{FAbsCCfbOne} percentage points more than the nominal fallback and more than a proportional buffer and a
safety factor tuned for this cost, and it improved moment propagation in the same way. At twice this cost, a tuned
buffer was better. Software-in-the-loop flights with PX4 and Gazebo confirmed the fallback, which raised the utility by
\val{HURelNoFb}\%, and agreed closely with the kinematic tier, but they did not resolve the advantage over the nominal
model on the missions flown. Future work will extend the cost-aware rule to the delay that a fallback task causes on
its route, fly it in the software-in-the-loop tier, run it inside a distributed allocator, estimate the disturbance
level online without loss of utility, and run the replanner on embedded hardware and in field flights.
